\chapter{Đặt vấn đề}

\section{Bối cảnh và động lực nghiên cứu}

Đi chung xe theo yêu cầu đã trở thành một trong những hướng được kỳ vọng nhiều nhất để giảm
áp lực giao thông đô thị. Ý tưởng rất trực tiếp: thay vì mỗi người một xe, hệ thống gom những
yêu cầu có lộ trình tương đồng vào cùng một phương tiện, nhờ đó giảm số xe lưu thông mà vẫn
phục vụ được cùng lượng nhu cầu. Công trình của Alonso-Mora và cộng sự \cite{alonsomora2017}
đã cho thấy tiềm năng của hướng này ở quy mô thật: với ba nghìn xe bốn chỗ, hệ thống của họ
phục vụ được khoảng $98\%$ nhu cầu taxi của Manhattan. Kết quả đó xác lập một chuẩn mực kỹ
thuật mà phần lớn các nghiên cứu sau này đều kế thừa, gồm đồ thị ba tầng Request--Trip--Vehicle
để sinh các khả năng ghép và một bài toán quy hoạch nguyên để chọn ra phương án tốt nhất.

Tuy nhiên, cách vận hành của các hệ thống này có một đặc điểm ít được nói tới. Chúng hoạt động
theo chu kỳ: cứ sau một khoảng thời gian ngắn, toàn bộ bài toán gán khách cho xe lại được giải
lại từ đầu với trạng thái mới. Mỗi lần giải lại, kế hoạch của một chiếc xe có thể thay đổi --
chèn thêm một điểm đón, đảo thứ tự hai điểm dừng, thậm chí chuyển một hành khách sang xe khác.
Điều đó là cần thiết để hệ thống tận dụng được mọi cơ hội ghép chuyến, và cũng chính là nguồn
gốc hiệu quả của nó.

Vấn đề nằm ở chỗ hành khách không nhìn thấy quá trình tối ưu ấy. Cái họ nhìn thấy là một con
số trên màn hình điện thoại: ``xe sẽ đón bạn lúc 9 giờ 10''. Khi hệ thống tối ưu lại và con số
đó trở thành 9 giờ 13, rồi 9 giờ 11, rồi 9 giờ 16, hành khách trải nghiệm một chuỗi lời hứa bị
rút lại chứ không phải một quá trình tối ưu. Và đây mới là điểm đáng chú ý: trong dòng bài điều
phối kinh điển, chuỗi ấy \emph{không được ghi nhận} trong mô hình. Các hệ thống này ràng buộc
hai tham số --- thời gian chờ tối đa $\Omega$ và độ trễ tối đa $\Delta$ --- và miễn là kế hoạch
còn nằm trong cửa sổ đó, nó được phép nhảy tùy ý, bao nhiêu lần cũng được. Một hành khách được
báo đúng một lần và một hành khách bị đổi giờ mười hai lần là hai đối tượng giống hệt nhau
trong mọi hàm mục tiêu của dòng bài này.

Sự thiếu vắng đó không phải chuyện nhỏ về mặt kinh tế. Một nhánh văn liệu khác, thuộc về phía
cầu, đã đo được rất rõ rằng hành khách trả giá cho độ tin cậy. Alonso-González và cộng sự
\cite{alonsogonzalez2020} ước lượng giá trị độ tin cậy cho dịch vụ đi chung theo yêu cầu;
Geržinič và cộng sự \cite{gerzinic2023} nghiên cứu cảm nhận của hành khách về độ biến thiên
thời gian chờ, và phần chờ phát sinh ngoài dự kiến bị đánh giá tiêu cực hơn phần chờ đã biết
trước.
Bansal và cộng sự \cite{bansal2019} cho thấy việc công bố ETA một cách trung thực làm tỉ lệ
chấp nhận chuyến tăng gần $10\%$. Nhưng chính các tác giả đó viết rõ trong bài rằng họ không
mô hình hóa chi phí phía cung: khung của họ giả định nhà điều hành \emph{luôn} có thể giao
khách đúng mốc đã hiển thị.

Nói cách khác, có cả một nhánh nghiên cứu chứng minh rằng giữ lời rất đáng tiền, và một nhánh
khác xây dựng những hệ thống điều phối hiệu quả; câu hỏi nối hai nhánh ấy lại mới chỉ bắt đầu
được trả lời gần đây (mục \ref{muc:chodung}):

\begin{khungy}
Giữ lời tốn bao nhiêu? Và liệu một lời hứa như thế có cưỡng chế được về mặt thuật toán hay
không, hay nó chỉ là một chỉ tiêu mong muốn?
\end{khungy}

\section{Hạn chế của các phương pháp hiện tại}

Để xác định chính xác chỗ đứng của khóa luận, chúng tôi đã khảo sát mười bốn trục văn liệu liên
quan. Từ đó rút ra bốn hạn chế mang tính hệ thống của dòng bài điều phối kinh điển, không phải
những thiếu sót ngẫu nhiên của một bài cụ thể; các ngoại lệ gần đây được bàn riêng ở mục
\ref{muc:chodung}.

\textbf{Thứ nhất, hành khách được mô hình hóa như một ràng buộc chứ không phải một đối tác.}
Cửa sổ $\Omega$ và $\Delta$ là tất cả những gì hệ thống nợ hành khách. Không tồn tại khái niệm
``điều đã nói với khách'' trong mô hình. Hệ quả trực tiếp là cửa sổ hợp lệ không hề bảo vệ lời
đã hứa: trong thực nghiệm của chúng tôi, một hành khách bị dời giờ trả tích lũy tới 573 giây
mà hệ thống vẫn tính là phục vụ thành công, vì mọi mốc vẫn nằm trong cửa sổ.

\textbf{Thứ hai, không có quy trách nguyên nhân.} Khi giờ hẹn của một hành khách xấu đi, không
ai tách phần do nhà điều hành quyết định lại khỏi phần do xe chạy chậm hoặc do đường tắc.
Fielbaum và Alonso-Mora \cite{fielbaum2020} là ngoại lệ đáng kể nhất --- họ phân loại nguyên
nhân thành ``do người dùng gây ra'' và ``do nhà điều hành gây ra'' --- nhưng đó là một phân
loại định tính bằng lời, và bản thân thước đo của họ thì loại tắc đường ra khỏi phạm vi ngay
từ đầu.

\textbf{Thứ ba, chỉ có thước đo kết quả, không có thước đo hành vi.} Bộ chỉ số chuẩn gồm tỉ lệ
phục vụ, thời gian chờ, quãng đường phương tiện và tỉ lệ ghép. Không chỉ số nào trả lời câu hỏi
``hệ thống này cư xử có giữ lời hay không''.

\textbf{Thứ tư, các khái niệm về ổn định lịch trình đã tồn tại ở ngành khác nhưng chưa được
chuyển sang lĩnh vực này.} Trong hoạch định sản xuất, hiện tượng ``MRP nervousness'' đã được
nghiên cứu từ thập niên 1980. Sridharan, Berry và Udayabhanu \cite{sridharan1987} đo độ ổn định
của lịch sản xuất chính dưới chân trời cuốn chiếu, và ``phương thuốc'' kinh điển của họ ---
đóng băng một phần lịch --- chính là cơ chế mà chúng tôi gọi là khóa lời hứa. Trong vận trù
giao hàng, bài toán Consistent VRP \cite{kovacs2014} có ràng buộc nhất quán giờ đến giữa các
ngày. Nhưng đó là bài toán tĩnh, nhiều ngày, cho hàng hóa; nó không nói gì về một lời hứa động
đưa ra cho một hành khách trong vòng vài phút.

\section{Chỗ đứng so với công trình gần nhất}
\label{muc:chodung}

Trước khi trình bày đóng góp, cần nói rõ ranh giới với ba công trình gần nhất, vì phần chồng lấn
là đáng kể và việc giấu nó đi sẽ chỉ làm người đọc tự phát hiện ra ở chương sau.

\paragraph{Fielbaum và Alonso-Mora (2020).} Fielbaum và Alonso-Mora \cite{fielbaum2020} đã đặt
đúng vấn đề nêu ở mục 1.1 và đã làm sáu việc
mà khóa luận này cũng làm: nêu khái niệm lời hứa công bố bị xê dịch, đo độ lớn dịch chuyển trên
từng hành khách, đếm số lần dự báo bị đổi, phân loại nguyên nhân thành do người dùng và do nhà
điều hành, dựng đường cong đánh đổi giữa mức kiểm soát và chất lượng dịch vụ, và thử cơ chế
ratchet trên giờ đón. Khóa luận coi sáu hạng mục ấy là \emph{nền}, và phần tương ứng trong
chương 4 là một cuộc \emph{tái lập độc lập} trên một thực thể khác --- có giá trị nhưng không
phải đóng góp mới.

Khác biệt nằm ở một điểm cấu trúc, và nó không phải chuyện mức độ mà là bản chất. Thước đo của
họ, công thức (2.5) ở chương sau, định nghĩa qua $t_w(r)$ --- thời gian chờ \emph{thực tế} ---
tức một đại lượng chỉ biết được sau khi hành khách đã được đón. Tại thời điểm hệ thống phải
chọn phương án nào, giá trị đó chưa tồn tại. Nghĩa là thước đo ấy, về mặt toán học, \emph{không
thể trở thành ràng buộc}: không ai chặn được ``thực tế trừ lần báo đầu không quá 30 giây'' khi
chưa biết thực tế. Nó chỉ tồn tại ở thì quá khứ.

Muốn chặn hiện tượng mà bài năm 2020 mô tả, đại lượng bị chặn phải định nghĩa trên những gì biết
được tại thời điểm quyết định. Riêng điều đó thì chưa cần gì mới: chính họ đã chặn theo cách ấy.
Kịch bản ``kiểm soát hoàn toàn'' ở mục 6 của bài đánh dấu khách đã được gán là không thể bị từ
chối và siết giờ chờ tối đa của họ xuống bằng giờ đã dự báo --- một đại lượng biết được lúc quyết
định; trên hệ 2.000 xe của họ, tỉ lệ từ chối tăng từ $34{,}5\%$ lên $45{,}6\%$, và mục ấy kết thúc
bằng nhận xét rằng cần những ý tưởng tinh vi hơn để kiểm soát bất ổn. Phép phân rã ba mốc ở
chương 3 tồn tại vì một lý do khác: tách phần do quyết định khỏi phần thay đổi khi giữ nguyên kế
hoạch, để một trần có thể chỉ tính phần thứ nhất.

\paragraph{Elmi, Moylan, Alonso-Mora và Fielbaum (2026).} Chính nhóm nghiên cứu ấy sau đó đưa ra
một câu trả lời \cite{elmi2026}, và đây là đối thủ gần nhất về ``bảo đảm lời hứa''. Họ dựng một
chỉ số ghép chuyến riêng cho từng hành khách để ước lượng khả năng người đó bị đổi giờ, rồi
\emph{cấm hoàn toàn} việc tăng thời gian đi cho nhóm có chỉ số thấp, còn nhóm có chỉ số cao giữ
nguyên độ chùng. Ở cùng tỉ lệ từ chối khoảng $29\%$, tỉ lệ khách bị đổi giờ giảm từ $63{,}17\%$
xuống $15{,}39\%$ --- khoảng 4,1 lần --- và $83{,}85\%$ khách được bảo đảm không đổi giờ. Nhưng mức
tăng trung bình của những người \emph{vẫn} bị đổi lại tăng từ $0{,}84$ lên $3{,}74$ phút ---
khoảng 4,5 lần: gánh nặng được \emph{phân phối lại} chứ không được xóa, điều chính các tác giả
ghi nhận như một vấn đề công bằng. Họ viết rằng một lời bảo đảm như vậy ``would not be possible if
even a single change were allowed''. Mạng của họ có thời gian đi trên mỗi cung cố định, nên không
có phần trôi dạt nào để tách khỏi phần do quyết định. Bài được nhận ngày 16/1/2026, chấp nhận ngày
24/7/2026 và công bố trực tuyến ngày 4/8/2026; hai công trình được thực hiện song song.

Một trần có mức cho mọi khách tự nó không phải điểm khác: ở Elmi và cộng sự, mọi lần chèn khách
mới đều không được vượt độ chùng còn lại của bất kỳ khách nào đã gán (mục 2.1.4 và 2.4.1 của bài),
và đối chứng ``ordinary control'' của họ cấp cho \emph{mọi} khách cùng một tỉ lệ độ chùng. Đó là
một trần trên phần tăng thời gian đi so với lúc gán, tức gần với họ $M^{+}$ ở mục \ref{muc:hocong}.
Khóa luận khác ở \emph{đối tượng} bị chặn: nó chặn \emph{tổng phần do nhà điều hành gây ra}, tách
khỏi phần thay đổi khi giữ nguyên kế hoạch. Trên mạng có thời gian đi cố định như của họ, hai đối
tượng gần như trùng nhau; dưới trôi dạt thì tách ra --- trần hành vi không chặn cái khách nhìn thấy
(phản ví dụ CE-2 ở bảng \ref{bang:phanvidu}), đổi lại nó không bao giờ tự nó làm rỗng tập phương
án. Với trần 30 giây, tổng gánh nặng hành vi giảm $99{,}8\%$
(từ $64.863$ xuống $128$ giây trên chiều giờ trả) thay vì được dồn sang một nhóm, với cái giá
$7{,}13$ điểm phần trăm tỉ lệ phục vụ. Hai con số ấy không so trực tiếp với số của Elmi và cộng
sự, vì hệ, dữ liệu và thước đo khác nhau; chỉ \emph{hình dạng} --- phân phối lại hay xóa --- là so
được.

\paragraph{Schulz và Pfeiffer (2026).} Đây là đối thủ gần nhất về ``trần dịch chuyển''
\cite{schulz2026}. Ngay khi khách chấp nhận, họ tính một hạn chót bằng giờ tới điểm trả đã báo
cộng một hằng số chung (mặc định 15 phút); mọi lần chèn khách mới về sau phải giữ giờ trả không
muộn hơn hạn chót ấy. Đó là một trần cứng trên phần \emph{muộn} hiện tại so với lời hứa \emph{đầu
tiên} (không cộng dồn các lần dời, và không chặn chiều sớm lên), kiểm trước khi chọn phương án --- tiền lệ mạnh nhất cho một ràng buộc lời hứa cưỡng chế được mà
chúng tôi tìm thấy. Trong mô hình của họ thời gian đi giữa hai điểm dừng là cố định, và bài khẳng
định khách đã nhận không bao giờ bị từ chối về sau.

\paragraph{Các công trình khác chạm tới lời hứa.} Haferkamp \cite{haferkamp2024} đưa ra nhiều lựa
chọn giờ đón và, giống Schulz và Pfeiffer, giới hạn mức giờ tới của khách đã nhận được phép trễ
đi; Sivagnanam và cộng sự \cite{sivagnanam2026} xác nhận nhanh các yêu cầu đặt trước và bảo đảm
yêu cầu đã nhận sẽ được phục vụ; Schmidt và cộng sự \cite{schmidt2026} chọn giờ đến công bố cho
khách trước khi khách quyết định đặt; Zalesak và cộng sự \cite{zalesak2025} nói về tính ổn định
của \emph{thuật toán} tính chi phí hoán đổi chứ không phải của lời hứa. Trong phần chúng tôi đã
đọc của các công trình này --- toàn văn với bài của Sivagnanam và của Zalesak, phần tóm tắt hoặc
phần được trích lại với hai bài còn lại --- không có sổ cái cộng dồn theo từng khách hay phép quy
trách định lượng nào.

\paragraph{Không gian thiết kế.} Đặt các cách làm cạnh nhau cho bốn điểm:
\begin{itemize}[leftmargin=1.2cm,itemsep=2pt,topsep=3pt]
\item cấm làm xấu giờ chờ đã báo cho mọi khách (Fielbaum và Alonso-Mora 2020, mục 6): bất ổn
về giờ chờ về không, tỉ lệ từ chối tăng $11{,}1$ điểm phần trăm trên hệ của họ;
\item cấm tăng cho một nhóm chọn lọc, giới hạn trong độ chùng cho những khách còn lại (Elmi và
cộng sự 2026): gánh nặng được phân phối lại; mạng của họ có thời gian đi cố định;
\item hạn chót tính từ lời hứa đầu (Schulz và Pfeiffer 2026): theo Định lý \ref{dl:hocong}, có
thể tự nó làm rỗng tập phương án khi thời gian đi trôi dạt theo chiều bất lợi --- một hệ quả của
định lý, chưa được quan sát trên hệ của họ vì mô hình của họ không có trôi dạt;
\item trần hành vi có mức cho mọi khách (khóa luận này): xét riêng phần cam kết, đây là cách duy
nhất trong bốn họ ràng buộc được so ở chương 3 \emph{không bao giờ là nguyên nhân} làm rỗng tập
phương án; đổi lại nó không chặn phần thay đổi khi giữ nguyên kế hoạch mà khách vẫn nhìn thấy.
\end{itemize}

\begin{khungy}
Chỗ khác biệt chứng minh được của khóa luận không nằm ở việc ``có bảo đảm lời hứa cho từng
khách'' --- Elmi và cộng sự đã có --- cũng không nằm ở việc chặn theo chuỗi thay vì theo hai đầu,
vì trên bộ sinh chèn thuần hai cách chặn gần như trùng nhau. Nó nằm ở chỗ đặt các họ ràng buộc lời
hứa --- theo hành vi, theo hạn chót, theo kết quả nhìn thấy --- vào cùng một khung, và chỉ ra rằng,
xét riêng phần cam kết, chỉ ràng buộc theo hành vi không bao giờ tự nó là nguyên nhân làm rỗng tập
phương án khi thế giới trôi dạt, còn họ nào trong các họ khác hỏng thì phụ thuộc \emph{chiều} trôi
(Định lý \ref{dl:hocong}).
\end{khungy}

Hai lưu ý về mức độ của phát biểu này. Chứng minh của định lý là sơ cấp; giá trị của nó nằm ở
khung so sánh và ở việc đối chiếu với hệ thật. Và phần đối chiếu là một thí nghiệm \emph{thăm dò có
tiền đăng ký} (mục \ref{muc:tang2}): dự đoán theo từng ô được niêm phong trước khi chạy các nhánh
được dự đoán, đúng ở $98/102$ cặp trên panel chính và ở $31/32$, $40/40$ cặp trên một tập giữ lại
chưa từng dùng để dựng quy tắc dự đoán; các lần chạy có trôi dạt không qua bộ kiểm chứng của dự án.

Từ khác biệt cấu trúc ấy nảy sinh một câu hỏi mà trước đó không có lý do gì để đặt ra:
\emph{khi biến thước đo thành ràng buộc, nên tính sổ theo cái gì?} Chương 3 chứng minh rằng câu
trả lời quyết định cổng cam kết có thể tự nó làm hệ thống kẹt hay không, và rằng cách tự nhiên
nhất để biến thước đo của văn liệu thành ràng buộc --- tính theo kết quả khách nhìn thấy --- lại
chính là một cách có thể kẹt.

\section{Mục tiêu và đóng góp của khóa luận}

Từ bốn hạn chế ở mục 1.2 và khác biệt cấu trúc ở mục 1.3, khóa luận đặt mục tiêu xây dựng và
đánh giá một cơ chế cho phép hệ thống điều phối \emph{đưa ra lời hứa mà nó cưỡng chế được},
đồng thời đo chính xác cái giá của việc đó. Cụ thể gồm ba mục tiêu.

Mục tiêu thứ nhất là thiết kế cơ chế \textbf{lời hứa có biên được chứng nhận}: một sổ cái cam
kết riêng cho từng hành khách, cộng dồn theo thời gian và đối chiếu với một trần $\beta$, được
kiểm tra ngay tại thời điểm ra quyết định chứ không phải sau sự việc. Mục tiêu thứ hai là xây
dựng phép \textbf{phân rã ba mốc} để tách phần dịch chuyển do nhà điều hành gây ra khỏi phần
thay đổi khi giữ nguyên kế hoạch, qua đó xác định chính xác cái gì nên bị tính vào sổ cái. Mục tiêu
thứ ba là đo bằng thực nghiệm có kiểm soát toàn bộ đường cong đánh đổi giữa mức độ chặt chẽ
của lời hứa và tỉ lệ phục vụ.

Các đóng góp chính của khóa luận gồm:

\begin{itemize}[leftmargin=1.2cm,itemsep=3pt,topsep=4pt]
\item \textbf{Ràng buộc cam kết cưỡng chế được.} Một sổ cái riêng cho từng hành khách, cộng dồn
theo chuỗi các lần công bố và đối chiếu với một trần $\beta$ ngay tại thời điểm quyết định.
Đại lượng bị chặn định nghĩa hoàn toàn trên những gì biết được trước khi chọn phương án, nên chặn
được, và trần áp cho \emph{mọi} khách chứ không chỉ một nhóm. Thực nghiệm cho thấy trần luôn được
tôn trọng ở mọi cấu hình --- $291 < 300$, $119 < 120$, $59 < 60$, $27{,}6 < 30$ giây.

\item \textbf{Khung bốn họ cổng và định lý về nguyên nhân làm rỗng tập phương án.} Chúng tôi đặt
bốn họ ràng buộc lời hứa --- theo hành vi, hạn chót một chiều kiểu Schulz và Pfeiffer, hạn chót
hai chiều, theo kết quả khách nhìn thấy --- vào cùng một khung, và chứng minh rằng, xét riêng phần
cam kết, chỉ cổng hành vi không bao giờ là nguyên nhân làm rỗng tập phương án khi thế giới trôi
dạt; họ nào trong ba họ còn lại hỏng phụ thuộc chiều trôi (Định lý \ref{dl:hocong}). Đi kèm là ba
định lý về hai cơ sở tính sổ cái: cơ sở hành vi không bao giờ bị ràng buộc cam kết làm rỗng, cơ sở
kết quả thì có thể, và khi một điều kiện (H) đúng ở mọi lần công bố thì chính sách tối ưu theo cơ
sở kết quả không phục vụ nhiều hơn. Điều kiện (H) \emph{không} được thỏa nguyên vẹn trên panel ---
nó hỏng ở 228 lần công bố đã chọn của B1 và 3 lần của C1 --- nên định lý thứ ba là một mệnh đề có
điều kiện, không phải một kết luận về panel. Trong thí nghiệm bơm tắc đường đo thật của
Manhattan (thăm dò), cơ sở kết quả vô nghiệm do cam kết ở 19 trên 20 ô, còn cơ sở hành vi không
vô nghiệm do cam kết ở ô nào --- dù 11 ô vẫn dừng vì một bất biến cửa sổ đón không thuộc cơ chế cam
kết. Một thí nghiệm thăm dò có tiền đăng ký kiểm thêm hai họ hạn chót: dưới tắc đường buổi sáng, hạn
chót 30 giây tính từ lời hứa đầu vô nghiệm do cam kết ở 19/20 ô; buổi tối, khi tắc đường giảm, hạn
chót một chiều vô nghiệm ở 0/8 ô còn hạn chót hai chiều ở 7/8 ô, đúng chiều định lý dự đoán. Mức 30
giây chặt hơn nhiều so với 15 phút mà Schulz và Pfeiffer dùng. Khi không có trôi dạt, hạn chót và trần
hành vi cùng 30 giây nhận đúng cùng một tập khách ở 20/20 ô, nên kết quả này nói về vùng lời hứa chặt
dưới tắc đường. Một phép quét thăm dò từ 30 giây tới 15 phút (mục \ref{muc:quet}) cho thấy hai luật
phục vụ gần như cùng số khách ở mọi mức; khác biệt nằm ở chỗ lời hứa có giữ được không. Ở mức chặt
dưới tắc đường tự tạo có cú sốc, hạn chót phá lời hứa của chính nó với $10$--$19\%$ khách. Trên tắc đường
đo thật, ở hai khung đã qua cao điểm, nó chỉ phá lời hứa ở mức 30 giây buổi sáng, với $2{,}2\%$ khách của ô ấy;
nhưng ở khung 7:00--9:00, khi tắc đường tăng nhanh nhất, nó phá lời hứa với $24\%$ khách ở 30 giây và $7{,}3\%$ ở
1 phút, trong khi hai luật vẫn chở đúng như nhau. Trần hành vi giữ lời hứa của nó theo cấu tạo. Chứng minh là
sơ cấp; bằng chứng thực nghiệm là thăm dò.

\item \textbf{Phép phân rã ba mốc và quy trách định lượng.} Bằng cách chiếu lại chính kế hoạch
cũ từ vị trí thật hiện tại của xe --- một phương án có thật trong tập lựa chọn --- hệ thống tách
được nguyên nhân của mỗi lần dời lời hứa theo từng quyết định và từng hành khách. Đây là điều kiện
kỹ thuật để hai đóng góp trên tồn tại, và tự nó cho một con số: trong chế độ không cam kết, phần
do bộ điều phối tự gây ra bằng $89{,}06\%$ lượng dịch chuyển mà hành khách nhìn thấy.

\item \textbf{Đường cong giá của cam kết.} Chín cấu hình cam kết khác nhau được đo trên cùng
một panel với cùng hạt giống, cho phép định lượng chi phí của từng mức chặt chẽ. Con số trung
tâm: một trần 30 giây cắt xáo trộn xuống còn $0{,}17\%$ mức ban đầu với giá $7{,}13$ điểm phần
trăm tỉ lệ phục vụ.

\item \textbf{Phân tích độ vững với cách đo.} Chín thước đo thuộc họ Sridharan (1988) được áp
dụng song song; thứ hạng các cấu hình không đổi dưới bất kỳ thước đo nào, trong khi độ lớn
chênh nhau tới hai bậc. Đi kèm là một cảnh báo phương pháp luận: tỉ lệ phục vụ nhiễu hơn nhiều
so với các thước đo ổn định, và một báo cáo chỉ dựa vào tỉ lệ phục vụ dễ đọc nhầm nhiễu thành
kết quả.

\item \textbf{Đối chiếu trực tiếp với Fielbaum và Alonso-Mora (2020).} Chúng tôi cài đặt lại
thước đo của họ trên dữ liệu của mình và chỉ ra rằng thước đo hai đầu của họ bỏ sót $49\%$ số
hành khách thực sự bị dời lời hứa dưới chế độ cam kết, đồng thời tái lập được kết luận định
tính trung tâm của họ trên một thực thể nhỏ hơn 185 lần.
\end{itemize}

Một điểm cần nói rõ ngay từ đầu, và sẽ được nhắc lại trong chương thực nghiệm: khóa luận này
\emph{không} tuyên bố vượt trội về tỉ lệ phục vụ. Mọi cấu hình có cam kết đều phục vụ ít hành
khách hơn cấu hình không cam kết. Đóng góp nằm ở chỗ khác --- ở việc biến một đại lượng trước
nay chỉ đo được sau sự việc thành một ràng buộc cưỡng chế được trước khi quyết định, và ở việc
chỉ ra loại ràng buộc lời hứa nào có thể tự nó làm hệ thống kẹt khi thế giới trôi dạt.

\section{Phạm vi và giới hạn}

Khóa luận giới hạn phạm vi như sau. Dữ liệu là bộ chuyến đi taxi công khai của Manhattan
(nguồn TLC, giấy phép CC-BY) đã được chuẩn hóa sẵn trong bộ FleetPy \cite{fleetpy2022}, khung
thời gian 08:00--10:00 của năm ngày liên tiếp từ 14 đến 18 tháng 11 năm 2018 (thứ Tư đến Chủ nhật,
tức gồm cả hai ngày cuối tuần). Quy mô mỗi ô thực
nghiệm là 108 yêu cầu, 8 xe sức chứa 4 chỗ, cửa sổ đón 600 giây, ràng buộc thời gian trên xe
bằng $1{,}5$ lần thời gian đi thẳng. Toàn bộ thực nghiệm dùng một hạt giống cố định nên hoàn
toàn tất định và tái lập được. Nhu cầu là dữ liệu thật, nhưng ``mạng'' của bộ dữ liệu thử không
phải mạng đường phố: nó là bao đóng thời gian đi giữa 96 điểm đã chọn, với tọa độ tổng hợp (mục
\ref{muc:hanche}).

Quy mô này nhỏ hơn đáng kể so với các thực nghiệm của Alonso-Mora và cộng sự. Đó là lựa chọn
có chủ ý: mục tiêu của khóa luận là so sánh có kiểm soát giữa các cấu hình cam kết trên cùng
một đầu vào, chứ không phải chứng minh khả năng mở rộng. Hệ quả là các con số tuyệt đối trong
báo cáo này không so sánh trực tiếp được với con số trong các bài quy mô lớn; chỉ có các so
sánh nội bộ giữa các cấu hình mới có ý nghĩa. Chương 4 sẽ trở lại điểm này khi đối chiếu với
văn liệu.

\section{Cấu trúc báo cáo}

Phần còn lại của báo cáo được tổ chức theo đúng mạch lập luận đã nêu ở mục \ref{muc:chodung},
tức đi từ chỗ
``thước đo không chặn được'' tới chỗ ``ràng buộc chặn được, và cái đó tốn bao nhiêu''.

Chương 2 trình bày cơ sở lý thuyết. Ngoài bài toán ghép chuyến và kiến trúc giải nó, chương này
dành một mục riêng để phân tích ba giới hạn của cách đo trong văn liệu, trong đó giới hạn thứ
hai --- thước đo chỉ tồn tại ở thì quá khứ --- là nền trực tiếp cho chương sau. Chương cũng vay
mượn hai khung khái niệm từ ngoài lĩnh vực: độ ổn định lịch trình trong hoạch định sản xuất, và
phân biệt nghĩa vụ phương tiện với nghĩa vụ kết quả trong luật hợp đồng.

Chương 3 là trọng tâm về phương pháp. Nó xây dựng phép phân rã ba mốc như lời giải cho vấn đề
``đại lượng bị chặn phải biết được tại thời điểm quyết định'', rồi từ đó dựng sổ cái cam kết,
chứng minh ba định lý về hai cơ sở tính sổ, rồi đặt chúng cùng các ràng buộc của văn liệu vào
khung bốn họ cổng. Chương kết thúc bằng kiến trúc cài đặt và cơ chế chứng cứ niêm phong.

Chương 4 trình bày thực nghiệm, tổ chức theo khẳng định chứ không theo thứ tự các đợt chạy;
bảng lộ trình ở đầu chương nối mỗi khẳng định với mục chứng minh nó. Chương gồm cả phần đối
chiếu ranh giới với công trình gần nhất, phần ghi nhận các hướng đã thử và phải từ bỏ, và phần
liệt kê hạn chế.

Chương 5 tổng kết đóng góp, nêu rõ những hạn chế còn lại và đề xuất các hướng phát triển.
